\chapter{Sequencing Space Programs}
\label{ch:26}

The orbital programs of the preceding chapters are not independent, and the order in which they are undertaken matters more than the decision to undertake any of them. The accelerator of Chapter 21 supplies the first part of the velocity that the skyhook of Chapter 23 completes. The skyhook delivers material to orbits from which a fabrication industry could build the structures of Chapter 24 and the collectors of Chapter 25. The collectors supply the energy for the fabrication that builds them, and the whole depends on a supply of material that cannot come from the Earth in the quantities required. This circularity is the central fact of the sequencing problem, and it is the reason that the answer cannot be a calendar. The sequence is determined by dependencies, by the cost and information that each stage yields, and by the institutions that must be in place before each stage can be undertaken legitimately. This chapter sets out the principles by which the sequence is to be chosen, and it closes the fifth part by connecting the orbital programs back to the standard of Chapter 16 and the constraint of present obligations.

The first principle is that present obligations are not postponed for the sake of orbital programs. A reader who has followed the earlier chapters will see why the principle needs to be stated. The scale of the orbital programs invites a form of reasoning in which the magnitude of the eventual benefit, energy and material in quantities that dwarf the economy of the Earth, is invoked to justify the diversion of resources from nearer needs. The reasoning is that of the longtermist argument which Chapter 31 examines in detail, and the principle that blocks it is the readiness rule of Chapter 16: no measure that satisfies a present obligation may depend on a program of the second or third readiness class. The orbital programs may therefore be funded, but from resources that do not compete with the floor of present obligations, and the plan must show that this is so. A corollary is that the programs earn continued support only by passing gates, and not by the size of their promised payoff.

The second principle is that the sequence is determined by reversibility and information. A stage that is cheap, reversible, and informative should precede a stage that is costly, irreversible, and no more informative. A short tether in low orbit that validates the dynamics of rotation and capture costs little relative to the programs it informs and can be abandoned. A pilot accelerator of modest speed, used to test the technology of the track and the handling of payloads, similarly costs a small fraction of the full installation, and its failure would teach what the design lacks. By contrast the dismantling of a planet is irreversible, and its justification cannot be settled by anything less than a long record of smaller operations in the same technologies. The structure that results is the stage-gate process familiar from the development of drugs and aircraft: stages are ordered so that the cheapest uncertainty-resolving steps come first, each stage ends in a gate at which an independent body decides whether to proceed, and each gate is defined by criteria stated in advance.

The calculation that underlies the stage-gate is backward induction, and its result is easily stated. At the final stage the decision is whether to spend the cost of the stage in return for the probability of success times the value of completion. At the preceding stage the value of completion is the value of the continuation, which is the maximum of zero and the value just computed, and so on back to the beginning. The value of a program with four stages with costs of one, three, ten, and thirty billion dollars, success probabilities of $0.5$, $0.4$, $0.5$, and $0.6$, and a final value of five hundred billion is $22.5$ billion when each stage is undertaken only if the preceding one has succeeded and the continuation value justifies the cost. The same program pre-funded, with the whole forty-four billion committed at the outset, has an expected value of $-14$ billion, because the later and larger costs are incurred even in the many cases in which an early stage fails. The difference of $36.5$ billion is the value of the option to stop, and it is the quantitative reason for gates. The calculation is also a test of honesty in proposals: a proposal that cannot specify the stages, the costs, and the probabilities on which it relies has not stated what it claims, and a proposal whose value is positive only if the probabilities are far higher than the evidence supports should be treated accordingly. The probabilities here are illustrative, and a real program would estimate them by reference to the history of comparable development efforts, adjusted by an independent review for the optimism that proposers display.

The third principle concerns the critical path. The dependencies among tasks form a directed graph, and the earliest completion of the program is the length of the longest path through it, since tasks off that path can be delayed up to their slack without delaying the whole. The implication is that effort should be concentrated on the critical path, and that a reduction in the duration of a task off the path does not shorten the program. A less obvious implication is that the critical path can change: when a technology advances, a task that had been fast becomes the bottleneck, and the plan must be revised. For the orbital programs the critical path is probably the production of material in space, since the other elements, from the accelerator to the collector, can be developed in parallel but all depend on it for their scale. The longest path is the bottleneck, and the most valuable early investment is the one that shortens it.

The fourth principle concerns authorization, and it is here that the argument of the book is most severely tested. The space beyond the Earth is governed by treaties of the 1960s and 1970s that were not drafted with these programs in view. The Outer Space Treaty of 1967~\cite{ost1967} prohibits the national appropriation of celestial bodies, requires states to bear international responsibility for the activities of their nationals, and prohibits the placing of weapons of mass destruction in orbit. The Liability Convention of 1972 makes states liable for damage caused by their space objects. The Moon Agreement of 1979, which would have declared the resources of the Moon the common heritage of mankind, has been joined by few states and by none of the principal space powers. The result is a regime that says clearly that no state may claim a planet and says much less clearly who may mine one, who is entitled to the energy of the Sun, and who may refuse. A program that proposed to disassemble Mercury, or to intercept a significant fraction of the Sun's output, would be an intervention in the common circumstances of all persons, and no state or group of states can authorize it by itself. The book's position, derived from the architecture of authority of Chapter 13, is that such a program requires an authorization that is multilateral, that provides for the participation of those not parties to it, that allows a party to refuse participation without loss of its existing rights, and that is revocable. There are precedents in the treaties governing Antarctica~\cite{antarctic1959} and the seabed, which have reserved areas of common interest to the benefit of all and subjected them to collective decision, though the compromises were negotiated when the resources in question were out of reach, and they would be tested severely if reach were extended.

The fifth principle concerns dual use, which was raised for the accelerator in Chapter 21 and applies to the others. A skyhook is a means of changing the orbits of objects at will; a swarm of collectors is a means of concentrating power; the capacity to move large masses in space is the capacity to direct them. The orbital programs create capabilities that would be of the highest strategic importance, and the possession of them by a single party would be a standing threat to others. The sequencing plan must accordingly include, at every gate, an assessment of the new capability's potential for misuse and a statement of the transparency and the verification regime proposed to address it, and the plan should be refused at any gate at which the regime has not been agreed.

The last principle is that the plan must be revisable. A sequence of decades that cannot be changed on the basis of what is learned is a prediction and not a plan. The program therefore requires, at each gate, a review that has the power to change the sequence, to defer later stages, and to reallocate resources, together with a standing process for recording the failures of the program, as Chapter 10 requires of any verification process. The documentation of a stage that failed is an asset of the program, and the institution that suppresses it harms the program more than the failure did.

\section*{Formalization}

Let a program have stages $i=1,\dots,n$, with cost $c_i>0$ incurred at the start of stage $i$, success probability $p_i\in(0,1]$ of stage $i$ given that it is undertaken, and terminal value $V>0$ on success of all stages. Define the continuation values $W_{n+1}=V$ and
\[
W_i=\max\{0,\,-c_i+p_iW_{i+1}\},\qquad i=n,n-1,\dots,1,
\]
interpreted as the value at the beginning of stage $i$, given that stages $1,\dots,i-1$ have succeeded and that the decision to undertake stage $i$ is still open.

\begin{proposition}
The policy "undertake stage $i$ if and only if $p_iW_{i+1}\ge c_i$" maximizes expected value among all policies that may stop at any gate. The value of the gated program is $W_1$. The value of the pre-funded program, in which all stages are paid for regardless of outcomes, is $\Pi=-\sum_ic_i+V\prod_ip_i$. Then $W_1\ge\max\{0,\Pi\}$, and $W_1-\Pi\ge0$ with equality only if every continuation decision is positive at every gate and the pre-funded program incurs no cost after a failure.
\end{proposition}

\begin{proof}
The recursion is Bellman's principle: the value of optimal play from stage $i$ is the better of stopping (value $0$) and undertaking the stage, which costs $c_i$ and with probability $p_i$ yields the optimal value from stage $i+1$. Induction from $n$ to $1$ establishes optimality. For the comparison, any pre-funded schedule is a policy that does not stop, and the gated policy is optimal among stopping policies, so $W_1\ge$ its value; and stopping at the start gives $0$. For the equality case, the gated policy that never stops after a failure incurs expected cost $\sum_ic_i\prod_{j<i}p_j$, which is less than $\sum_ic_i$ whenever some $p_j<1$, $j<n$, and $n\ge2$.
\end{proof}

For $c=(1,3,10,30)$, $p=(0.5,0.4,0.5,0.6)$, and $V=500$ (in billions), $W_5=500$, $W_4=\max\{0,-30+0.6\cdot500\}=270$, $W_3=-10+0.5\cdot270=125$, $W_2=-3+0.4\cdot125=47$, and $W_1=-1+0.5\cdot47=22.5$. The pre-funded value is $\Pi=-44+500\cdot0.06=-14$. The expected expenditure of the gated program is $1+0.5\cdot3+0.2\cdot10+0.1\cdot30=7.5$ billion, against $44$ billion for the pre-funded one, and its expected value is $0.06\cdot500-7.5=22.5$ billion, in agreement. The gated program is positive at every gate here; if the estimate of $p_3$ were revised at the third gate from $0.5$ to $0.03$ the continuation value would fall to $-10+0.03\cdot270=-1.9$ and the policy would stop, avoiding the $40$ billion of the remaining outlay in an event the pre-funded program cannot respond to.

For the critical path, let the tasks form a directed acyclic graph $G=(T,E)$ with durations $d_t\ge0$. The earliest finish time $F_t$ of task $t$ satisfies $F_t=d_t+\max_{s\to t}F_s$, with the maximum over the empty set equal to $0$.

\begin{proposition}
The earliest completion time of the program is $\max_tF_t$, which equals the maximum over all paths of the sum of durations along the path. A task has \emph{slack} $\ell_t=\mathrm{LF}_t-F_t\ge0$ (latest finish minus earliest finish), and increasing $d_t$ by $\delta\le\ell_t$ leaves the completion time unchanged, while reducing $d_t$ for a task of positive slack does not reduce it. A reduction of $\delta$ of a task on every longest path reduces the completion time by at most $\delta$, and by exactly $\delta$ if the next-longest path is at least $\delta$ shorter.
\end{proposition}

The proof is the standard one for the critical path method and is omitted. An illustration: let the accelerator pilot (8 years) and the tether demonstration (6 years) both precede in-space fabrication (10 years), which precedes the collector pilot (12 years). The longest path is accelerator pilot, fabrication, pilot, of length $30$ years, and the tether demonstration has slack $2$ years, so a delay of up to two years in the tether program does not delay the whole, whereas a one-year gain in the accelerator pilot shortens the program by one year, up to the point where the tether path (28 years) becomes critical.
